% Generated draft: requires mathematical and source review.
\documentclass[11pt]{article}
\usepackage[margin=25mm]{geometry}
\usepackage{fontspec}
\setmainfont{Latin Modern Roman}
\usepackage{amsmath,amssymb,mathtools,unicode-math}
\setmathfont{Latin Modern Math}
\usepackage{enumitem,needspace,xurl,hyperref}
\setlength{\parindent}{0pt}
\setlength{\parskip}{5pt}
\setlength{\emergencystretch}{3em}
\newcommand{\N}{\mathbb{N}}
\newcommand{\Z}{\mathbb{Z}}
\newcommand{\Q}{\mathbb{Q}}
\newcommand{\R}{\mathbb{R}}
\newcommand{\C}{\mathbb{C}}
\newcommand{\F}{\mathbb{F}}
\newcommand{\A}{\mathbb{A}}
\newcommand{\PP}{\mathbb{P}}
\newcommand{\E}{\mathbb{E}}
\newcommand{\eps}{\varepsilon}
\newcommand{\dd}{\,\mathrm{d}}
\begin{document}
\section*{UnsolvedMath ID: 30001645}
\textbf{Source problem code:} catalogue entry
\subsection{Definitions and mathematical statement}
Let \(K/\Q\) be a finite abelian extension and put \(G=\operatorname{Gal}(K/\Q)\). Let \(\mu_K=\{\zeta\in K^\times:\zeta^n=1\text{ for some }n\geq1\}\) and \(e=|\mu_K|\). Write \(\operatorname{Cl}_K\) for the group of nonzero fractional ideals of the ring of integers of \(K\), modulo principal fractional ideals. The natural action of \(G\) on ideals makes \(\operatorname{Cl}_K\) a \(\Z[G]\)-module.

The supplied record does not explicitly specify the underlying field of the places in \(S\), the finiteness of \(S\), or the Frobenius convention. To give a precise candidate formulation, use the conventional interpretation in which \(S\) is a finite set of places of \(\Q\) containing the infinite place and every finite prime ramified in \(K/\Q\). For a finite prime \(p\notin S\), let \(\sigma_p\in G\) be arithmetic Frobenius, characterized on residue fields by \(x\mapsto x^p\). Since \(G\) is abelian, this element is independent of the prime of \(K\) chosen above \(p\). Here the norm of the base-field prime is \(Np=p\). These conventions require confirmation against the source; they are not additional hypotheses explicitly supplied by the record.

For \(\operatorname{Re}(s)>1\), define the equivariant Euler product
\[
\theta_S(s)=\prod_{p\notin S}\left(1-\sigma_p^{-1}p^{-s}\right)^{-1}\in\C[G],
\]
where the product runs over finite primes of \(\Q\) outside \(S\). Interpret \(\theta_S(0)\) as the value at zero obtained by analytic continuation, rather than by direct substitution in the Euler product.

\textbf{Annihilation question under these conventions.} For every finite abelian extension \(K/\Q\) and every such finite set \(S\), is
\[
 e\theta_S(0)\in\operatorname{Ann}_{\Z[G]}(\operatorname{Cl}_K)?
\]
Here
\[
\operatorname{Ann}_{\Z[G]}(\operatorname{Cl}_K)
=\{\alpha\in\Z[G]:\alpha\cdot c=0\text{ for every }c\in\operatorname{Cl}_K\},
\]
using additive notation for the class group. Equivalently, integral group-ring annihilation means that \(e\theta_S(0)=\sum_{\sigma\in G}n_\sigma\sigma\) with \(n_\sigma\in\Z\), and that for every nonzero fractional ideal \(\mathfrak a\) of \(K\), the fractional ideal
\[
\prod_{\sigma\in G}\sigma(\mathfrak a)^{n_\sigma}
\]
is principal. Integrality is implicit in this interpretation of annihilation; the supplied statement does not discuss it separately.

The record calls this the Brumer--Stark conjecture, but supplies only the class-group annihilation assertion. No further condition on generators of principal ideals is specified, so no stronger conjectural formulation is included here.
\subsection{Short English statement}
For a finite abelian extension \(K/\Q\), the supplied question asks whether multiplying the value at zero of an equivariant Euler product by the number of roots of unity in \(K\) gives an integral group-ring element annihilating every ideal class. The precise candidate uses a finite set of base-field places and arithmetic Frobenius; these conventions need source confirmation. The supplied partially solved label and dated literature assessment are curation claims, not independently verified conclusions.
\subsection{Sources}
\begin{itemize}
\item Supplied local UnsolvedMath registry record, permanent ID \(30001645\), exact source problem code \texttt{catalogue entry}, title: Brumer--Stark Class Group Annihilation Conjecture. No named proposer is supplied.
\item Supplied report citation: \emph{Mini-Workshop: 1-Motives (2011)}, Oberwolfach Reports, pp.~3237--3267, DOI \texttt{10.4171/owr/2010/55}. \url{https://doi.org/10.4171/owr/2010/55}. The background describes Oberwolfach Reports as documenting mathematical workshops and their open-problem discussions. The same URL is supplied under the literature-source label ``original Oberwolfach problem report.''
\item Supplied literature source, labeled ``Brumer-Stark away from 2'': \url{https://annals.math.princeton.edu/2023/197-1/p05}.
\item Literature assessment dated exactly \texttt{2026-08-21}: the supplied curation labels the problem partially solved and states that the Brumer--Stark conjecture is proved for the minus part away from \(2\), while the full integral conjecture in the stated generality remains unresolved. This is supplied dated curation derived from the OWR material, not an independently verified assessment. Its time-sensitive status, cited claims, and applicability to the displayed annihilation question require independent verification.
\end{itemize}
\end{document}
